% ECONOMICS_PROBLEM: {"id": "D91-587", "title": "Context dependence of social preferences", "jel_code": "D91", "source_id": "OP-0064"}
% !TeX program = xelatex
\documentclass[11pt,a4paper]{article}
\usepackage[margin=23mm]{geometry}
\usepackage{fontspec}
\setmainfont{Latin Modern Roman}
\usepackage{amsmath,amssymb,mathtools}
\usepackage{xcolor,enumitem,booktabs,longtable,array,xurl,hyperref}
\definecolor{muted}{HTML}{53636C}
\hypersetup{colorlinks=true,urlcolor=blue,linkcolor=blue}
\setlength{\parindent}{0pt}
\setlength{\parskip}{5pt}
\newcommand{\R}{\mathbb R}
\newcommand{\E}{\mathbb E}
\newcommand{\N}{\mathbb N}
\newcommand{\ind}{\mathbf 1}
\newcommand{\Var}{\operatorname{Var}}
\newcommand{\Cov}{\operatorname{Cov}}
\newcommand{\supp}{\operatorname{supp}}
\DeclareMathOperator*{\argmin}{arg\,min}
\DeclareMathOperator*{\argmax}{arg\,max}
\newcommand{\entrymeta}[3]{{\sffamily\small\color{muted}\textbf{\texttt{#1}}\quad|\quad #2\par #3\par}}
\newcommand{\Problemref}[1]{\texttt{#1}}
\newcommand{\JELref}[2]{\texttt{#2} (JEL \texttt{#1})}
\begin{document}
\section*{Context dependence of social preferences}
\textbf{Problem ID:} \texttt{D91-587}\quad \textbf{JEL:} \texttt{D91}\par
% BEGIN ECONOMICS STATEMENT
\label{problem:OP-0064}
\label{atlas:prob:B4}

\noindent\textbf{Original identifier:} B4 (atlas item 64). \textbf{Classification:} Social preference identification research program. \textbf{Original tractability label:} Medium (editorial, not a complexity result).

\subsection*{Definitions and assumptions}

Let $G$ specify an extensive-form game, monetary payoff vector $\pi$, information, repetition, anonymity, and framing. For $n\geq2$ players, a benchmark inequity-aversion representation is
\[
U_i(\pi)=\pi_i-\frac{a_i}{n-1}\sum_{j\ne i}(\pi_j-\pi_i)_+
-\frac{b_i}{n-1}\sum_{j\ne i}(\pi_i-\pi_j)_+,
\]
where $(z)_+=\max\{z,0\}$, $a_i\geq b_i\geq0$, and $b_i<1$. Reciprocity requires additional, explicitly specified beliefs about intentions; altruism permits utility from others' payoffs without inequity penalties. Let $\theta_i$ collect the declared social-preference parameters. A response model maps utilities and beliefs to observed actions, incorporating prespecified decision noise.

\subsection*{Formal research question}

For a declared collection $\mathcal G$ of laboratory and real-world interactions, determine whether one context-invariant latent preference distribution $F_\theta$ and correctly specified context-dependent beliefs can reproduce all action distributions:
\[
\forall G\in\mathcal G:\quad P_G(A\in E)=\int P_G(A\in E\mid\theta)\,dF_\theta(\theta)
\]
for every measurable action event $E$. Identify $F_\theta$ or characterize its observationally equivalent set. Compare the common-parameter model with a prespecified alternative $F_{\theta\mid G}$, using out-of-sample predictive loss and equivalence tolerances rather than unrestricted contextual fitting.

\subsection*{Interpretation and scope}

An action can change with context even when preferences remain fixed: beliefs about reputation, future interaction, or another person's intentions may change. Conversely, stable dictator allocations need not establish stable reciprocity. Games must therefore vary payoff inequality, efficiency, intention, and strategic incentives separately wherever feasible. Repeated within-person designs require controls for learning and carryover, and cross-cultural comparisons need common measurement and comprehensible instructions. The inequity-aversion formula is only a benchmark submodel; it is not a mathematical definition covering fairness, altruism, and reciprocity simultaneously. A successful empirical program would identify which contextual changes shift beliefs and which require preference variation, with uncertainty from both stages. Where models cannot be distinguished, reporting an identified set is a substantive conclusion. Recommendations for organizational design additionally require a specified welfare aggregator; willingness to punish is not itself proof that punishment improves welfare.

\subsection*{Source audit and status}

The three author-year entries resolve to inequity aversion [S1], experimental social-preference discrimination [S2], and altruistic punishment [S3]. Numbered link [38] is a different Fehr--Gächter paper, the 2000 reciprocity review [S4], and is corrected explicitly. These sources provide models and evidence, not a universal stable social-preference theorem. Context-invariant identification across the chosen game collection remains an editorial research objective whose status must be established for each application.

\subsection*{References}

\begin{enumerate}[label={[S\arabic*]},leftmargin=*]

\item Ernst Fehr and Klaus M. Schmidt (1999). \emph{A Theory of Fairness, Competition, and Cooperation}. Quarterly Journal of Economics 114(3), 817--868. \url{https://doi.org/10.1162/003355399556151}\par \textit{Source role:} Inequity-aversion model. \textit{Locator:} Abstract and article metadata.

\item Gary Charness and Matthew Rabin (2002). \emph{Understanding Social Preferences with Simple Tests}. Quarterly Journal of Economics 117(3), 817--869. \url{https://doi.org/10.1162/003355302760193904}\par \textit{Source role:} Experimental discrimination among social-preference models. \textit{Locator:} Abstract and article metadata.

\item Ernst Fehr and Simon Gächter (2002). \emph{Altruistic Punishment in Humans}. Nature 415, 137--140. \url{https://doi.org/10.1038/415137a}\par \textit{Source role:} Evidence on punishment and cooperation. \textit{Locator:} Abstract and article metadata.

\item Ernst Fehr and Simon Gächter (2000). \emph{Fairness and Retaliation: The Economics of Reciprocity}. Journal of Economic Perspectives 14(3), 159--181. \url{https://www.aeaweb.org/articles?id=10.1257/jep.14.3.159}\par \textit{Source role:} Original numbered link [38]; reciprocity review. \textit{Locator:} Abstract and article metadata.

\end{enumerate}

\subsection*{Common conventions: Source mathematical conventions}

Unless an entry states otherwise, agent and firm sets are finite. State and action spaces are measurable, strategies and outcome maps are measurable, probability laws are specified, and expectations exist for the targets under discussion. Infinite-horizon discount factors lie in $(0,1)$ and discounted objectives are finite. Norms, tolerances and complexity input sizes must be specified for computational comparisons. Constants or primitives that appear locally are inputs, not empirical facts asserted by this catalogue.

\textbf{Identification.} A statistical model $m\in\mathcal M$ implies an observable law $P_m$ and target $\theta(m)$. At law $P$, its identified set is
\[
\Theta(P;\mathcal M)=\{\theta(m):m\in\mathcal M,\ P_m=P\}.
\]
Point identification holds when this set is a singleton, or equivalently when $P_{m_1}=P_{m_2}$ implies $\theta(m_1)=\theta(m_2)$ for all compatible models. Sharp bounds mean bounds attained, or approached when the boundary is not attained, by compatible models. Writing this set is a definition, not a solution: the substantive task is to establish informative restrictions and derive or estimate the set. An unrestricted-confounding model may give only trivial bounds.

\textbf{Inference.} Identification, estimability and valid uncertainty quantification are separate questions. Fix $\alpha\in(0,1)$. A confidence procedure $C_n$ over a declared law class $\mathcal P$ has asymptotic uniform coverage at level $1-\alpha$ when
\[
\liminf_{n\to\infty}\inf_{P\in\mathcal P}\Pr_P\{\theta(P)\in C_n\}\geq1-\alpha.
\]
For set-identified targets the coverage event must be defined separately, for example coverage of the full identified set. If an entry uses identification-indexed models instead of uniquely indexed laws, the same coverage convention is applied over those models. Pointwise limits do not establish uniform validity, and asymptotic validity does not guarantee finite-sample precision. Sample sizes, dependence structures and weak-identification sequences are part of each application.

\textbf{Counterfactuals.} $Y_i(a)$ is a potential outcome under a specified intervention, including its timing, implementation and versions. It is not $Y_i$ conditional on voluntarily choosing $a$. A full intervention vector permits interference; shorthand scalar treatment notation does not silently impose no interference. Assignment, selection, attrition, anticipation and measurement must be specified. A claim about an instrument's local effect is not automatically a population policy effect.

\textbf{Economic equilibrium.} A counterfactual model includes best responses, constraints, feasibility, market clearing, state transitions and expectations consistent with the stated information. When several equilibria exist, either a declared selector is an input or the target is set-valued. Comparisons across policy regimes must use the stated selection convention. Reduced-form relationships are not presumed invariant after an equilibrium-changing intervention.

\textbf{Optimization and welfare.} Optimization is over the stated feasible set. If attainment is not established, $\sup$ or $\inf$ and $\varepsilon$-optimal choices replace an asserted maximizer or minimizer. For example, with value $V=\sup_{a\in\mathcal A}W(a)<\infty$, an $\varepsilon$-optimal set is $\{a\in\mathcal A:W(a)\geq V-\varepsilon\}$ for $\varepsilon>0$. Benefits, costs, risk and health or environmental outcomes can be aggregated only using declared units and conversion weights. Interpersonal utility weights, the treatment of fiscal transfers and foreign profits, discounting, and the behavioral welfare criterion are explicit inputs. A comparative-static derivative additionally requires existence and appropriate differentiability of the selected counterfactual.

\textbf{Sources.} Local labels [S1], [S2], and so forth restart in every question. Locators identify the relevant abstract, model, section or result at the level actually checked; a bibliographic or abstract-level verification is not represented as inspection of an entire proof. Working papers are distinguished from later publications and changing versions. Source summaries are paraphrased. All URL checks and editorial formulations refer to the audit date stated above.
% END ECONOMICS STATEMENT
\end{document}
